\documentclass{article}
\usepackage[T1]{fontenc}
\usepackage{lmodern}
\usepackage{graphicx}
\usepackage{amsmath,amssymb}
\usepackage[a4paper,margin=2cm]{geometry}
\usepackage[absolute,overlay]{textpos}
\setlength{\TPHorizModule}{1mm}
\setlength{\TPVertModule}{1mm}
\usepackage[indonesian]{babel}
\usepackage{hyperref}
\setlength{\emergencystretch}{2em}
% unit: Halaman 28

\begin{document}

% === Halaman 28 (python/native) ===

28

Bentuk umum cipher abjad-majemuk

• Kunci:        

              K = k1k2 … km                       (ket: m = panjang kunci)

   Karena m lebih pendek daripada panjang plainteks, maka kunci diulang secara 

periodik

• Plainteks:  

              P = p1p2 … pmpm+1 … p2m …                       

• Cipherteks:

          𝐶= 𝐸𝑘𝑃= 𝑓𝑘1(𝑝1) 𝑓𝑘2(𝑝2)… 𝑓𝑘𝑚(𝑝𝑚)𝑓𝑘1(𝑝𝑚+1)𝑓𝑘2(𝑝𝑚+2)… 𝑓𝑘𝑚(𝑝2𝑚)…


f(.) adalah fungsi enkripsi pada cipher abjad-tunggal  


• Untuk m = 1, cipher-nya ekivalen dengan cipher abjad-tunggal.

\end{document}
